\documentclass[10pt]{extarticle}
\usepackage[a4paper,margin=1.5cm]{geometry}
\usepackage[T1]{fontenc}
\usepackage{amsmath,amssymb}
\usepackage{graphicx}
\usepackage[export]{adjustbox}
\usepackage{multicol}
\usepackage{enumitem}
\usepackage{booktabs,array,tabularx}
\setlength{\columnsep}{0.7cm}
\setlength{\parindent}{0pt}
\setlength{\parskip}{3pt}
\pagestyle{empty}
\makeatletter
% figure registry (figdefs.tex is generated by the extract step)
\newcommand{\qfigdef}[4]{\@namedef{qf@file@#1}{#2}\@namedef{qf@w@#1}{#3}\@namedef{qf@wide@#1}{#4}}
\newcommand{\qf@check}[1]{\@ifundefined{qf@file@#1}{\PackageError{score2}{Unknown figure id #1}{}}{}}
\newcommand{\qf@img}[4][]{\includegraphics[width=#2\dimexpr\@nameuse{qf@w@#4}pt\relax,max width=#3,max height=0.45\textheight,#1]{\@nameuse{qf@file@#4}}}
% Full-width questions: finalize wraps every question that contains a wide figure in \qfullbegin...\qfullend,
% so the WHOLE question (stem, figures, options) is set across the page in reading order, then the two columns
% resume. The question's head - its number, stem and first figure (\qfullsplit marks the end of the head) -
% is kept on one page, so a number is never left at a page foot with its figure overleaf; the rest may break
% across pages normally, so no big gaps are left. Without a \qfullsplit, the whole question is kept together
% if it is under 45% of a page. Never write these yourself.
\newif\ifqf@full
\newif\ifqf@boxed
\newsavebox{\qf@box}
\newcommand{\qf@open}{\global\qf@fulltrue\par\global\qf@boxedtrue
  \setbox\qf@box\vbox\bgroup\hsize\textwidth\linewidth\textwidth}
\newcommand{\qf@emit}[1]{\par\egroup\global\qf@boxedfalse
  \ifdim\dimexpr\ht\qf@box+\dp\qf@box\relax<#1\textheight\noindent\usebox{\qf@box}\par\else\unvbox\qf@box\fi}
\newcommand{\qfullsplit}{\ifqf@boxed\qf@emit{0.85}\fi}
\newcommand{\qf@close}{\ifqf@boxed\qf@emit{0.45}\fi\global\qf@fullfalse}
\newcommand{\qfullbegin}{\end{multicols}\qf@open}
\newcommand{\qfullend}{\qf@close\begin{multicols}{2}}
\newcommand{\qfullstart}{\qf@open}
\newcommand{\qfulldone}{\qf@close}
% \qfig[scale]{id}: block figure, centred, directly where it is written (size set by the tool).
% Fallback if a wide figure is not inside a full-width question: break out of the columns in place - no floats.
\newcommand{\qfig}[2][1]{\qf@check{#2}%
  \ifqf@full\qf@blk{#1}{#2}\else
  \ifnum\@nameuse{qf@wide@#2}=1 \expandafter\@firstoftwo\else\expandafter\@secondoftwo\fi
  {\end{multicols}\begin{center}\qf@img{#1}{\textwidth}{#2}\end{center}\begin{multicols}{2}}%
  {\qf@blk{#1}{#2}}\fi}
\newcommand{\qf@blk}[2]{\par\smallskip{\centering\qf@img{#1}{\linewidth}{#2}\par}\smallskip}
% \qfigi[scale]{id}: inline figure (inside option lists / tables / minipages); never breaks columns.
\newcommand{\qfigi}[2][1]{\qf@check{#2}\qf@img[valign=c]{#1}{\linewidth}{#2}}
\makeatother
% structure macros - lint reads these to check completeness and figure placement
\newcommand{\papersec}[1]{\par\bigskip{\centering\Large\bfseries PAPER-#1\par}\medskip}
\newcommand{\partsec}[2]{\par\medskip\noindent{\large\bfseries PART-#1: #2}\par\smallskip}
\newcommand{\sect}[1]{\par\smallskip\noindent\textbf{#1}\par}
\newcommand{\qn}[1]{\par\smallskip\noindent\textbf{#1.}~}
\newcommand{\paraq}[2]{\par\smallskip\noindent\textbf{Paragraph for Questions #1 and #2}\par}
\newenvironment{qoptions}{\begin{itemize}[leftmargin=1.8em,itemsep=1pt,topsep=2pt,label={}]}{\end{itemize}}
\newcommand{\opt}[2]{\item[(#1)] #2}
% figure registry (generated by score2_tools.py finalize)
\qfigdef{p1_phy_s0_q03_1}{SRGBT - Test 08 (13 Sep 2021)-img/p1_phy_s0_q03_1.png}{76.1}{0}
\qfigdef{p1_phy_s0_q05_1}{SRGBT - Test 08 (13 Sep 2021)-img/p1_phy_s0_q05_1.png}{196.0}{0}
\qfigdef{p1_phy_s0_q06_1}{SRGBT - Test 08 (13 Sep 2021)-img/p1_phy_s0_q06_1.png}{75.3}{0}
\qfigdef{p1_phy_s0_q06_2}{SRGBT - Test 08 (13 Sep 2021)-img/p1_phy_s0_q06_2.png}{321.5}{1}
\qfigdef{p1_phy_s0_q07_1}{SRGBT - Test 08 (13 Sep 2021)-img/p1_phy_s0_q07_1.png}{118.3}{0}
\qfigdef{p1_phy_s0_q08_4}{SRGBT - Test 08 (13 Sep 2021)-img/p1_phy_s0_q08_4.png}{306.3}{1}
\qfigdef{p1_phy_s0_q10_1}{SRGBT - Test 08 (13 Sep 2021)-img/p1_phy_s0_q10_1.png}{92.9}{0}
\qfigdef{p1_phy_s0_q11_1}{SRGBT - Test 08 (13 Sep 2021)-img/p1_phy_s0_q11_1.png}{79.8}{0}
\qfigdef{p1_phy_s0_q12_1}{SRGBT - Test 08 (13 Sep 2021)-img/p1_phy_s0_q12_1.png}{101.2}{0}
\qfigdef{p1_phy_s0_q14_1}{SRGBT - Test 08 (13 Sep 2021)-img/p1_phy_s0_q14_1.png}{150.7}{0}
\qfigdef{p1_phy_s0_q15_1}{SRGBT - Test 08 (13 Sep 2021)-img/p1_phy_s0_q15_1.png}{141.4}{0}
\qfigdef{p1_phy_s0_q17_1}{SRGBT - Test 08 (13 Sep 2021)-img/p1_phy_s0_q17_1.png}{94.6}{0}
\qfigdef{p1_chem_s0_q19_1}{SRGBT - Test 08 (13 Sep 2021)-img/p1_chem_s0_q19_1.png}{109.3}{0}
\raggedcolumns
\makeatletter\@beginparpenalty=10000\makeatother
\newcolumntype{L}[1]{>{\raggedright\arraybackslash}p{#1}}
\begin{document}
\begin{center}{\Large\bfseries SRGBT --- Test 08 (13 Sep 2021)}\end{center}
\begin{multicols}{2}
\raggedcolumns
\papersec{1}
\partsec{1}{PHYSICS}
\qn{1} A particle is projected with an initial velocity of u at an angle $\theta$ w.r.t. the horizontal. The tangential acceleration of the particle after time t is given by :
\begin{qoptions}
\opt{A}{$\mathrm g\dfrac{\mathrm u\cos\theta}{\sqrt{\mathrm u^2+\mathrm g^2\mathrm t^2-2\mathrm{ugt}\sin\theta}}$}
\opt{B}{$\mathrm g\dfrac{(\mathrm u\sin\theta-\mathrm{gt})}{\sqrt{\mathrm u^2+\mathrm g^2\mathrm t^2-2\mathrm{ugt}\sin\theta}}$}
\opt{C}{$\mathrm g\dfrac{\mathrm u\cos\theta}{\sqrt{\mathrm u^2+\mathrm g^2\mathrm t^2-2\mathrm{ugt}\cos\theta}}$}
\opt{D}{$\mathrm g\dfrac{(\mathrm u\sin\theta-\mathrm{gt})}{\sqrt{\mathrm u^2+\mathrm g^2\mathrm t^2-2\mathrm{ugt}\cos\theta}}$}
\end{qoptions}
\qn{2} A swimmer crosses a river with minimum possible time 10 second. And when he reaches the other end starts swimming in the direction towards the point from where he started swimming. Keeping the direction fixed the swimmer crosses the river in 15 sec. The ratio of speed of swimmer with respect to water and the speed of river flow is (Assume constant speed of river \& swimmer):-
\begin{qoptions}
\opt{A}{$\dfrac32$}
\opt{B}{$\dfrac94$}
\opt{C}{$\dfrac2{\sqrt5}$}
\opt{D}{$\dfrac{\sqrt5}2$}
\end{qoptions}
\qn{3} A particle is projected from point A at an angle of 53$^\circ$ with horizontal. At the same time wedge starts from rest and moves with constant acceleration a as shown. The value of a for which the particle strikes the plane perpendicular to it at P, is (tan53$^\circ$ = 4/3, g = 10 m/s$^2$)\par
\qfig{p1_phy_s0_q03_1}
\begin{qoptions}
\opt{A}{10 m/s$^2$}
\opt{B}{100 m/s$^2$}
\opt{C}{$\dfrac{130}{15}$ m/s$^2$}
\opt{D}{$\dfrac{140}{15}$ m/s$^2$}
\end{qoptions}
\qn{4} If speed (V), acceleration (A) and force (F) are considered as fundamental unit's, the dimension of pressure will be--
\begin{qoptions}
\opt{A}{V$^{-2}$A$^2$F$^2$}
\opt{B}{V$^{-4}$A$^2$F}
\opt{C}{V$^{-4}$A$^{-2}$F}
\opt{D}{V$^{-2}$A$^2$F$^{-2}$}
\end{qoptions}
\qn{5} A particle is projected from origin with speed u=50 m/s at an angle 30$^\circ$ with X-axis. Due to wind particle acquires an extra acceleration a = 6 m/s$^2$ in $-$ve x-direction. Path of particle is shown in figure.\par
\qfig{p1_phy_s0_q05_1}
x - Coordinates of point A and B are same\par
y- Coordinates of P and Q are same.\par
Point N on the X-axis, point M on the Y-axis.\par
Time taken by particle from O to P, O to Q, O to A and O to B are t$_1$, t$_2$, t$_3$ and t$_4$ respectively. (Assuming Y along vertical and their is no surface in the space) (g = 10 m/s$^2$ and constant).
\begin{qoptions}
\opt{A}{The value of t$_1$ + t$_2$ is 5 sec.}
\opt{B}{The value of t$_3$ + t$_4$ is 25 sec.}
\opt{C}{The length of ON is $125\sqrt3\,(5-\sqrt3)$ m}
\opt{D}{The length of OM is $\dfrac{625}{\sqrt3}\left(\dfrac5{\sqrt3}-1\right)$ m}
\end{qoptions}
\qfullbegin
\qn{6} Particle A and B are moving with constant velocities along positive x and positive y axis respectively, shown in figure. The graph of separation and the relative velocity between them with time is represented by (Assume that they will not collide):\par
\qfig{p1_phy_s0_q06_1}
\qfullsplit
\qfig{p1_phy_s0_q06_2}
\qfullend
\qn{7} When a rectangular box is going down a smooth inclined plane with initial velocity v$_\mathrm o$, a person standing on the rectangular box throws up a ball perpendicular to the inclined plane with velocity u with respect to himself as shown in the figure. Which of the following is/are true about the ball, coming back to incline?\par
\qfig{p1_phy_s0_q07_1}
\begin{qoptions}
\opt{A}{It will fall behind the person}
\opt{B}{It will fall ahead of the person}
\opt{C}{It will come back into the hands of person}
\opt{D}{Neglecting the height of point of projection above box total time taken by the ball to return on the box is $\dfrac{2\mathrm u}{\mathrm g\cos\theta}$}
\end{qoptions}
\qfullbegin
\qn{8} A ball is projected from ground with speed u, at an angle $\theta$ above horizontal. Let v be its speed at any moment `t' and `s' be the total distance covered by it till this moment, the correct graph(s) is / are :\par
\qfig{p1_phy_s0_q08_4}
\qfullend
\qn{9} Two particles A and B are projected from the same point with the same velocity of projection but at different angles $\alpha$ and $\beta$ of projection respectively, such that the maximum height of A is two-third of the horizontal range of B. Then which of the following relations is\par
Choose correct statement(s) :
\begin{qoptions}
\opt{A}{3(1 $-$ cos 2$\alpha$) = 8 sin 2$\beta$}
\opt{B}{Range of A = maximum height of B}
\opt{C}{Maximum value of $\beta$ is $\dfrac12\sin^{-1}\dfrac34$}
\opt{D}{Maximum horizontal range of A = $\dfrac{\mathrm u^2}{\mathrm g}$}
\end{qoptions}
\qn{10} A smooth inclined plane has an angle of inclination $\theta$ from the horizontal. From point A on the inclined plane, body 1 is projected horizontally with speed u and from point B on the inclined plane, body 2 is released from rest. Distance AB is equal to $\ell$. If the magnitude of relative velocity of body 2 w.r.t. body 1 and magnitude of relative acceleration of body 2 w.r.t. body 1 at time t are v and `a', respectively, then which of the following is/are true?\par
\qfig{p1_phy_s0_q10_1}
\begin{qoptions}
\opt{A}{v = $\sqrt{\mathrm u^2-2\mathrm{ug}\sin\theta\cos\theta\mathrm t+\mathrm g^2\cos^2\theta\,\mathrm t^2}$}
\opt{B}{a = zero}
\opt{C}{a = g cos $\theta$}
\opt{D}{v = $\sqrt{(\mathrm g\sin^2\theta\,\mathrm t-\mathrm u)^2+(\mathrm g\sin\theta\cos\theta\,\mathrm t-\mathrm{gt})^2}$}
\end{qoptions}
\qn{11} A shell of mass 1 kg is projected with velocity 20 ms$^{-1}$ at an angle 60$^\circ$ with horizontal. It collides inelastically with a ball of mass 1 kg which is suspended through a thread of length 1 m. The other end of the thread is attached to the ceiling of a trolley of mass 18 kg as shown in the figure. Initially, the trolley is stationary and it is free to move long horizontal rails without any friction. The maximum deflection of the thread with vertical is $\theta$ such that $\cos\theta=\dfrac1\alpha$, Taking, g = 10 ms$^{-2}$, find the value of $\alpha$\par
\qfig{p1_phy_s0_q11_1}
\qn{12} A boat travels across a river from a point A to a point B of opposite bank along the line AB forming angle 60$^\circ$ with the bank. A flag on the pole (mounted on the boat) flutters at an angle 60$^\circ$ with the line AB as shown in figure. If velocity of wind ($5\sqrt3$ m / s ) is perpendicular to the flow of river, find the speed (in m/s) of boat with respect to the bank.\par
\qfig{p1_phy_s0_q12_1}
\qn{13} A particle is projected from ground towards a vertical wall 80 m away at an angle of 37$^\circ$ with horizontal with initial velocity of 50 m/s. After its collision with wall \& then once with ground find at what distance d in meter from wall will it strike the ground. The component of velocity normal to the surface becomes half after collision with each surface n = $\dfrac{\mathrm d}{35}$. Find n.
\qn{14} In the given figure points A and C are on the horizontal ground \& A and B are in same vertical plane. Simultaneously bullets are fired from A, B and C and they collide at D. The bullet at B is fired horizontally with speed of $\dfrac{72}5$ km/hr and the bullet at C is projected vertically upward at velocity of $\dfrac{54}5$ km/hr. Find velocity of the bullet projected from A in m/s.\par
\qfig{p1_phy_s0_q14_1}
\qn{15} In the vertical plane shown two particles `P' and `Q' are located at points `A' and `B'. At t = 0, the particle `P' is projected perpendicular to the inclined plane `OA' with velocity 2v$_0$ and simultaneously the particle Q is projected horizontally in `BO' direction. the necessary value of v$_0$ (in terms of `d' and $\theta$) so that both the particles meet each other between the points `O' and `B' is $\mathrm v_0=\sqrt{\dfrac{\mathbf X\mathrm{gd}}{(1+2\sin\theta)[8\cos\theta+\tan\theta(1+2\sin\theta)]}}$ then calculate value of X.\par
\qfig{p1_phy_s0_q15_1}
\qn{16} Three particles A, B and C are situated at the vertices of an equilateral triangle. They start moving with equal speeds (remains constant during the whole motion) such that A always heads towards B, B towards C \& C towards A. Finally they meet at a point. Find out the angle by which the line connecting A and B rotaes with respect to initial orientation in space by the time distance between A and B reduces to half of the initial distance between A and B is $\pi/\alpha$. Find $\alpha$.
\qn{17} Find the speed of the intersection point O of the two wires if the wires starts moving perpendicular to itself with speed v as shown in figure is v cosec ($\alpha$/x). Find x.\par
\qfig{p1_phy_s0_q17_1}
\qn{18} $\displaystyle\int\dfrac{\mathrm{dv}}{\sqrt{2\mathrm{av}-\mathrm v^2}}=\mathrm a^\mathrm n\sin^{-1}\left[\dfrac{\mathrm x}{\mathrm a}-1\right]$ on the basis of dimensional analysis, the value of n is
\raggedcolumns
\papersec{1}
\partsec{2}{CHEMISTRY}
\qn{19} To 500 ml H$_3$PO$_4$ solution, 1 M Ca(OH)$_2$ solution was slowly added and the following plot was obtained. The initial molarity of H$_3$PO$_4$ solution and volume of Ca(OH)$_2$ required for complete neutralization are respectively. (Assuming complete ionization of H$_3$PO$_4$)\par
\qfig{p1_chem_s0_q19_1}
\begin{qoptions}
\opt{A}{$\dfrac23$,500 m$\ell$}
\opt{B}{$\dfrac23$,1500 m$\ell$}
\opt{C}{$\dfrac43$,2000 m$\ell$}
\opt{D}{$\dfrac13$,500 m$\ell$}
\end{qoptions}
\qn{20} In Raymonds factory, H$_2$O$_2$ is used for bleaching clothes. For avoiding errors, bottles are prepared according to the type of cloth. Each bottle has volume 500 ml and only one bottle is used for 1000 gm cloth. Information about bottles is given in the following chart :\par
{\small\setlength{\tabcolsep}{4pt}\noindent\begin{tabular}{|p{1.5cm}|p{2.7cm}|p{2.7cm}|}\hline Type of cloth & Required mass of O$_2$ per 100 gm of cloth & Name of Bottle to bleach 1 kg cloth\\\hline Cotton & 3.2 gm & P\\\hline Polyster & 4.8 gm & Q\\\hline Silk & 5.6 gm & R\\\hline\end{tabular}}\par
If H$_2$O$_2$ is produced by given reaction,\par
O$_3$ $\xrightarrow{h\nu}$ O$_2$ + O \hfill \dots(i)\par
H$_2$O$_2$ + O $\longrightarrow$ H$_2$O$_2$ \hfill \dots(ii)\par
then calculate the number of photons required in the first reaction to produce sufficient amount of nascent oxygen [O] to prepare one bottle of ``\textbf{R}'' type, if quantum yield of the first reaction is 0.2. [Given : N$_\mathrm A$ = 6 $\times$ 10$^{23}$]
\begin{qoptions}
\opt{A}{2.1 $\times$ 10$^{25}$}
\opt{B}{1.05 $\times$ 10$^{24}$}
\opt{C}{2.1 $\times$ 10$^{24}$}
\opt{D}{1.05 $\times$ 10$^{25}$}
\end{qoptions}
\qn{21} For a redox reaction, FeS$_2$ + KMnO$_4$ $\longrightarrow$ Fe$_2$O$_3$ + S$_8$ + Mn$^{2+}$ if 1 mole of S$_8$ is produced then calculate moles of KMnO$_4$ consumed
\begin{qoptions}
\opt{A}{1 mole}
\opt{B}{$\dfrac85$ mole}
\opt{C}{$\dfrac{16}5$ mole}
\opt{D}{$\dfrac{12}5$ mole}
\end{qoptions}
\qn{22} In a compound A$_\mathrm x$B$_\mathrm y$, correct statement is:
\begin{qoptions}
\opt{A}{Mole of A = mole of B = mole of A$_\mathrm x$B$_\mathrm y$.}
\opt{B}{Eq. of A = eq. of B = eq. of A$_\mathrm x$B$_\mathrm y$.}
\opt{C}{y $\times$ mole of A = x $\times$ mole of B = (x + y) $\times$ mole of A$_\mathrm x$B$_\mathrm y$.}
\opt{D}{y $\times$ mole of A = y $\times$ mole of B.}
\end{qoptions}
\qn{23} Given following series of reactions:\par
(I) NH$_3$ + O$_2$ $\longrightarrow$ NO + H$_2$O\par
(II) NO + O$_2$ $\longrightarrow$ NO$_2$\par
(III) NO$_2$ + H$_2$O $\longrightarrow$ HNO$_3$ + HNO$_2$\par
(IV) HNO$_2$ $\longrightarrow$ HNO$_3$ + NO + H$_2$O\par
Select the correct option(s):
\begin{qoptions}
\opt{A}{Moles of HNO$_3$ obtained is half of moles of Ammonia used if HNO$_2$ is not used to produce HNO$_3$ by reaction (IV)}
\opt{B}{$\dfrac{100}6$\% more HNO$_3$ will be produced if HNO$_2$ is used to produce HNO$_3$ by reaction (IV) than if HNO$_2$ is not used to produce HNO$_3$ by reaction (IV)}
\opt{C}{If HNO$_2$ is used to produce HNO$_3$ then $\dfrac14$th of total HNO$_3$ is produced by reaction (IV)}
\opt{D}{Moles of NO produced in reaction (IV) is 50\% of moles of total HNO$_3$ produced.}
\end{qoptions}
\qn{24} Which of the following solutions will be definitely basic.
\begin{qoptions}
\opt{A}{400 ml of 1.0 M NaOH is mixed with 100 ml of 0.3 M CH$_3$COOH}
\opt{B}{Equal volumes of same\% w/v of NaOH and H$_2$SO$_4$ solutions are mixed.}
\opt{C}{Different volumes of 0.2 M NaOH and 0.1 M HCl solution are mixed.}
\opt{D}{Equal masses of same \% w/w NaOH and HCl solutions are mixed.}
\end{qoptions}
\qn{25} Consider two reaction : (X represents an element)\par
\textbf{Reaction - I :} \quad H$_2$XO$_4$ + KOH $\to$ K$_2$XO$_4$ + H$_2$O\par
\textbf{Reaction - II :} \quad H$_2$XO$_4$ + KI $\to$ H$_2$X + I$_2$\par
Which of following statement(s) is/are incorrect -
\begin{qoptions}
\opt{A}{Molar mass of H$_2$XO$_4$ is equal in both reactions}
\opt{B}{Equivalent mass of H$_2$XO$_4$ is equal in both reactions}
\opt{C}{Equivalent mass of H$_2$XO$_4$ in Reaction I is twice as in Reaction -II}
\opt{D}{One mole of H$_2$XO$_4$ will contain same number of equivalent in both reactions}
\end{qoptions}
\qn{26} Choose incorrect statement(s) from the following :
\begin{qoptions}
\opt{A}{K$_4$[Fe(CN)$_6$] $\longrightarrow$ K$^+$ + Fe$^{+++}$ + CO$_2$ + NO$_3^-$ , equivalent weight of K$_4$[Fe(CN)$_6$] is $\dfrac{\mathrm M}{61}$.}
\opt{B}{Fe$_2$(C$_2$O$_4$)$_3$ $\longrightarrow$ Fe$^{+++}$ + CO$_2$ , equivalent weight of Fe$_2$(C$_2$O$_4$)$_3$ is $\dfrac{\mathrm M}7$}
\opt{C}{Na$_2$S$_2$O$_3$ $\longrightarrow$ Na$_2$SO$_4$ , equivalent weight of Na$_2$S$_2$O$_3$ is $\dfrac{\mathrm M}2$.}
\opt{D}{SnCl$_2$ + HCl + O$_3$ $\longrightarrow$ SnCl$_4$ + H$_2$O , equivalent weight of O$_3$ is $\dfrac{\mathrm M}2$}
\end{qoptions}
\qn{27} `4V' mL of 1M Na$_2$SO$_4$ is mixed with `V' mL of 4M Ba(NO$_3$)$_2$ solution. Select the incorrect statement(s).
\begin{qoptions}
\opt{A}{Molarity of Ba$^{2+}$ \& NO$_3^-$ in final solution are $\dfrac45$M \& $\dfrac85$M respectively.}
\opt{B}{Molarity of NO$_3^-$ \& Na$^+$ in final solution are $\dfrac85$M \& $\dfrac85$M respectively.}
\opt{C}{Molarity of Na$^+$ ion in final solution can't be calculated as V is not known.}
\opt{D}{Molarity of SO$_4^{2-}$ \& NO$_3^-$ in final solution are $\dfrac85$M \& $\dfrac85$M respectively.}
\end{qoptions}
\qn{28} In 20 ml of a gaseous mixture containing N$_2$ and H$_2$ gases, 5 ml O$_2$ gas is added and the mixture is exploded. If the final volume becomes 13ml, then the correct statement is/are :\par
(All the volumes are measured at the same pressure and temperature)
\begin{qoptions}
\opt{A}{The initial mixture contains 10ml N$_2$ gas.}
\opt{B}{The initial mixture contains 8ml H$_2$ gas.}
\opt{C}{The final mixture contains 1ml O$_2$ gas.}
\opt{D}{The final mixture contains 12ml N$_2$ gas.}
\end{qoptions}
\qn{29} In 20 ml of solution of NaBrO$_3$, 168 ml of SO$_2$ gas (measured at 0$^\circ$C and 1 atm) is passed, which reduced some NaBrO$_3$ to NaBr. The resulting solution required 50 ml of 0.2 N solution of a reducing agent for complete reduction or remaining NaBrO$_3$ into NaBrO. The concentration (in g/L) of NaBrO$_3$ solution taken was\par
[Br = 81, Na = 23, O = 16]
\qn{30} Beautician of `X' have 4.45 $\ell$ H$_2$O$_2$ solution of molality `\textbf{2m}' \& density 1.2 gm/m$\ell$. 2.5 mol of O$_2$ is used for bleaching of `X' hair from the solution.\par
Calculate Q, where Q = 89 $\times$ volume strength of the remaining H$_2$O$_2$ solution.\par
(Assume amount of water produced to be negligible and there is no change in density of the solution)
\qn{31} If Cl$_2$(g) and PCl$_3$(g) can react to form PCl$_5$ according to following reaction.\par
Cl$_2$(g) + PCl$_3$(g) $\longrightarrow$ PCl$_5$(g)\par
In tube A, 355 gm Cl$_2$ is taken and in tube B, 1375 gm PCl$_3$ is taken. When half the mass of gas in tube A is transferred to B, reaction took place in B. After the reaction is complete, half of the mass of mix. in B is again transferred to A to cause same reaction in container A. (no reversibility).\par
The mass fraction of PCl$_5$ in tube B.
\qn{32} Once Tom and Jerry entered a chemistry lab where a chemist was preparing 3 litre H$_2$SO$_4$ solution. He labelled the solution as \textbf{10 m (d$_{\textbf{Solution}}$ = 3.3 gm/ml)}. As the chemist left the lab, a mischief came in Tom's mind. He tired to throw the solution on Jerry but failed. In doing so some of the H$_2$SO$_4$ solution fell on the floor, so he added water to make it again to 3 litre. The chemist returned back \& got astonished when he saw the result of analysis, that were \textbf{d = 1.5 gm/ml} and \textbf{\% w/w = 98}. Find out the moles of H$_2$SO$_4$ which fell down on the floor.
\qn{33} Urnaium is isolated from its ore by dissolving it as UO$_2$(NO$_3$)$_2$ and separating it as solid UO$_2$(C$_2$O$_4$).x H$_2$O.A 1.0 g sample of ore on treatment with nitric acid yielded 1.48 g UO$_2$(NO$_3$)$_4$ which on further treatment with 0.4 g Na$_2$C$_2$O$_4$ yielded 1.23 g UO$_2$(C$_2$O$_4$).x H$_2$O. Determine weight percentage of uranium in the original sample and x.
\qn{34} In one process for waterproofing, a fabric is exposed to (CH$_3$)$_2$SiCl$_2$ vapour. The vapour reacts with hydroxyl groups on the surface of the fabric or with traces of water to form the waterproofing film [(CH$_3$)$_2$SiO]$_\mathrm n$, by the reaction\par
n(CH$_3$)$_2$SiCl$_2$ + 2nOH $\longrightarrow$ 2nCl$^-$ + nH$_2$O + [(CH$_3$)$_2$SiO]$_\mathrm n$\par
where a stands for a large integer. The waterproofing film is deposited on the fabric layer upon layer. Each layer is 6.0 \AA\ thick [the thickness of the (CH$_3$)$_2$SiO group]. How much (CH$_3$)$_2$SiCl$_2$ is needed to waterproof one side of a piece of fabric, 1.00 m by 3.00 m, with a film 300 layers thick ? The density of the film is 1.0 g/cm$^3$.\par
\textbf{First multiply your answer with 100. Then fill your answer as sum of digits (excluding decimal places) till you get the single digit answer.}
\qn{35} A polystyrene, having the formula Br$_3$C$_6$H$_3$(C$_3$H$_8$)$_\mathrm n$, was prepared by heating styrene with tribromobenzoyl peroxide in the absene of air. If it was found to contain 10.46\% bromine by weight, find the value of n.
\qn{36} An experiment was carried out to determine the hardness of domestic water. A few drops of pH = 10 buffer were added to a 100 cm$^3$ sample of the water. The sample was titrated against 0.01 M EDTA using Eriochrome Black T indicator and gave a titre of 33.8 cm$^3$. Under the circumstances both Mg$^{2+}$ and Ca$^{2+}$ ions react with EDTA. A second 100 cm$^3$ sample was titrated against EDTA solution after 5 cm$^3$ of 0.1 M NaOH and a few drops of murexide indicator had been added. Under these circumstances only the Ca$^{2+}$ ions react with EDTA and a titre of 27.5 cm$^3$ was obtained. Determine the hardness of water.
\raggedcolumns
\papersec{1}
\partsec{3}{MATHEMATICS}
\qn{37} The number of ordered pairs (x, y) satisfying the system of equations $6^x\left(\dfrac23\right)^y-3(2^{x+y})-8(3^{x-y})+24=0$; xy = 2 is:
\begin{qoptions}
\opt{A}{0}
\opt{B}{1}
\opt{C}{4}
\opt{D}{3}
\end{qoptions}
\qn{38} If equation $|x^2-5x+6|-\lambda x+7\lambda=0$ has exactly 3 solutions then $\lambda$ is equal to :
\begin{qoptions}
\opt{A}{$-7+\sqrt{23}$}
\opt{B}{$-9+4\sqrt5$}
\opt{C}{$-7-\sqrt{23}$}
\opt{D}{$-9-4\sqrt5$}
\end{qoptions}
\qn{39} Number of real values of x satisfying the equation $25^{\log_{125}(2\sqrt2)}+\log_{|x|}\left(\dfrac{|x|-\sqrt3}{\sqrt3|x|-2\sqrt2}\right)=\log_5\left(\dfrac{\sqrt3|x|+2\sqrt2}{|x|+\sqrt3}\right)\log_{|x|}5$ is:
\begin{qoptions}
\opt{A}{2}
\opt{B}{4}
\opt{C}{6}
\opt{D}{8}
\end{qoptions}
\qn{40} (6!)! is divisible by (6!)$^{\mathrm P!}$ then the greatest value of P is
\begin{qoptions}
\opt{A}{5}
\opt{B}{6}
\opt{C}{8}
\opt{D}{2}
\end{qoptions}
\qn{41} Which of the following statement is correct. If S = 1!+2!+3!+4!+\dots..1997! then :
\begin{qoptions}
\opt{A}{Unit digit of S is 3}
\opt{B}{Unit digit of S is 2}
\opt{C}{Remainder obtained when S is divide by 24 is 9}
\opt{D}{Remainder obtained when $2^{100}+3^{100}+4^{100}+5^{100}$ divided by 7 is 5}
\end{qoptions}
\qn{42} Considering the equation $\dfrac2{\mathrm x}+\dfrac5{\mathrm y}=\dfrac13$ where x, y $\in$ N. The number of solution of the equation is :
\begin{qoptions}
\opt{A}{6, if both x and y are even}
\opt{B}{6, if both x and y are odd}
\opt{C}{0 if x is even, and y is odd}
\opt{D}{0 if x is odd, and y is even}
\end{qoptions}
\qn{43} The equation $|x+1|\,|x-1|=a^2-2a-3$ can have real solution in x if a belongs to :
\begin{qoptions}
\opt{A}{$[1-\sqrt5,-1]$}
\opt{B}{$[-1, 3]$}
\opt{C}{$[3, 1+\sqrt5]$}
\opt{D}{$[1+\sqrt3, 1+\sqrt5]$}
\end{qoptions}
\qn{44} The equation $x^{\frac34(\log_2x)^2+\log_2x-\frac54}=\sqrt2$ has :
\begin{qoptions}
\opt{A}{at least one real root}
\opt{B}{exactly three real solutions}
\opt{C}{exactly one irrational solution}
\opt{D}{exactly one integral solution}
\end{qoptions}
\qn{45} Which of the following statement is correct (given log 2 = 0.3010, log 3 = 0.4771)
\begin{qoptions}
\opt{A}{number of digits in $18^{20}$ is 26}
\opt{B}{number of zeroes between decimal and first significant digit in $\left(\dfrac16\right)^{100}$ is 77}
\opt{C}{Number of values of x satisfying the equation $\log_4(x+3)-\log_4(x-1)=\log_4 8-2$ is 1}
\opt{D}{The value of $\log_{1/3}\sqrt[4]{729\sqrt[3]{9^{-1}\times27^{-4/3}}}$ is $-1$}
\end{qoptions}
\qn{46} If a, b, c be the sides of a triangle and equation $x^2+y^2-2x-4y-4+(a+b+c)\left(\dfrac1a+\dfrac1b+\dfrac1c\right)=0$ in variables x and y has real solution then :
\begin{qoptions}
\opt{A}{the triangle is equilateral}
\opt{B}{the triangle is of constant area}
\opt{C}{exactly one ordered pair (x, y) of solution is possible}
\opt{D}{the equation can never have real solution}
\end{qoptions}
\qn{47} The number of negative integers belonging to the solution set of the equation $2^{|x+2|}-|2^{x+1}-1|=2^{x+1}+1$ is equal to
\qn{48} The least positive integer x, whcih satisfies the inequality $\log_{\log_2\left(\frac{\mathrm x}2\right)}(x^2-10x+22)>0$ is equal to :
\qn{49} Number of integral values of x satisfying the equation $\sqrt{x+3-4\sqrt{x-1}}+\sqrt{x+8-6\sqrt{x-1}}=1$ is
\qn{50} The number of values of real x satsifying $\log_{3x+4}(4x^2+4x+1)+\log_{2x+1}(6x^2+11x+4)=4$
\qn{51} $\mathrm S=\dfrac{\sqrt1}{1+\sqrt1+\sqrt2}+\dfrac{\sqrt2}{1+\sqrt2+\sqrt3}+\dfrac{\sqrt3}{1+\sqrt3+\sqrt4}+\dots..$n terms = 3 then the values of n(n $\in$ N) is :
\qn{52} The value of n (n $\in$ N) for which $3^9+3^{12}+3^{15}+3^\mathrm n$ is a perfect cube is \underline{\hspace{2cm}}
\qn{53} If xy = $2^2\times3^4\times5^7$(x+y) the number of integral solutions is \underline{\hspace{2cm}}:
\qn{54} Minimum number of squares that can be cut from a rectangular board of dimension 130 $\times$ 50 such that no part of board is wasted is :
\end{multicols}
\end{document}
